\documentclass[10pt,a4paper]{amsart}
\usepackage[margin=19mm]{geometry}
\usepackage{amsmath,amssymb,mathtools}
\usepackage[scr=rsfs,cal=euler]{mathalfa}
\usepackage{xfrac}
\usepackage[unicode,hidelinks,bookmarks=false]{hyperref}
\newcommand{\ie}{i.e.\ }
\newcommand{\cf}{cf.\ }
\newcommand{\resp}{respectively\ }
\newcommand{\Subsection}{\subsection}
\providecommand{\nobreakdash}{\nobreak\hbox{-}}
\setlength{\emergencystretch}{2em}
\multlinegap0pt
\usepackage{amssymb}
\usepackage{bm}
\usepackage{mathtools}
\usepackage{tcolorbox}
\usepackage{subfig}
\usepackage[all]{xy}
\usepackage{bbm}
\usepackage{tikz}
\usepackage{cancel}
\usepackage{float}
\usepackage{xcolor}
\usepackage{multirow}
\usepackage{booktabs}
\usepackage[shortlabels]{enumitem}
\newtheorem{lem}{Lemma}
\DeclareMathOperator{\id}{\mathrm{id}}
\title{Non-isomorphic metacyclic $p$-groups of split type with the same group zeta function}
\author{Yuto Nogata}
\date{Source: arXiv:2512.24546v1 (CC BY 4.0)}
\begin{document}
\maketitle

\noindent\emph{Source and licence.} Non-isomorphic metacyclic $p$-groups of split type with the same group zeta function. Version arXiv:2512.24546v1 (CC BY 4.0), distributed by arXiv under the Creative Commons Attribution 4.0 International licence, http://creativecommons.org/licenses/by/4.0/, as recorded in the arXiv licence metadata for this exact version and shown on the abstract page https://arxiv.org/abs/2512.24546v1. Full licence notice: \texttt{licenses/arxiv-2512.24546-CC-BY-4.0.txt}.\par\smallskip

\noindent\textit{Editorial scope.} Counted unit: the article's lem lem:p-group-lattice-iso-implies-zeta (source0.tex lines 1205--1212). The article's complete original proof of it is delivered with it (source0.tex lines 1214--1234, one \texttt{proof} environment of 21 lines). No mathematics is written, repaired or re-derived: the statement and the proof are the article's own lines, in the article's own notation, and no retained line is substituted or re-flowed.  No further source-local context is needed: every cross-reference of every retained line resolves inside the two delivered spans.  Nothing was omitted from any retained span, so the package carries no omission record.  No retained line contains a citation, so the wrapper carries no bibliography and no bibliography entry.  No retained line contains a figure environment, an image-inclusion command or drawing code, so no image is delivered and the package declares no asset dependency.  Editorial formatting only, in addition: an AMS \texttt{amsart} preamble that restates the article's own declarations and macros used by the retained lines, namely the environments \texttt{lem} on separate counters, together with the standard packages those lines need.  The retained environments print in this document as Lemma 1 in order of appearance, whereas the article prints the same items with the numbering of its own full document.  The complete native source of the article as distributed in the arXiv e-print for this exact version is bundled unchanged as \texttt{sources/191907/source0.tex}.
\begin{lem}\label{lem:p-group-lattice-iso-implies-zeta}
Let $p$ be a prime and let $G_{1}$ and $G_{2}$ be finite $p$-groups.
Assume that $|G_{1}|=|G_{2}|$.
Then the following implication holds.
\[
  L(G_{1})\simeq L(G_{2}) \Longrightarrow \zeta_{G_{1}}(s) = \zeta_{G_{2}}(s).
\]
\end{lem}

\begin{proof}
Let $G$ be a finite $p$-group, and let $H\subset K\subset G$.
Then $H\prec K$ in $L(G)$ is equivalent to saying that $H$ is a maximal subgroup of $K$.
Since every maximal subgroup of a finite $p$-group has index $p$, it follows that
$H\prec K$ if and only if $[K:H]=p$.

Let $G_{1}$ and $G_{2}$ be finite $p$-groups with $|G_{1}|=|G_{2}|$,
and let $\Phi\colon L(G_{1})\to L(G_{2})$ be a lattice isomorphism.
Let $A\subset G_{1}$, and take a maximal chain $\{\id\}=A_{0}\prec A_{1}\prec \dotsb \prec A_{r}=A$
in the interval $[\{\id\},A]$.
Then $[A_{i}:A_{i-1}]=p$ for each $i$, hence $|A|=p^{r}$ and $\log_{p}|A|=r$.
Since $\Phi$ preserves cover relations, the chain
$\{\id\}=\Phi(A_{0})\prec \Phi(A_{1})\prec \dotsb \prec \Phi(A_{r})=\Phi(A)$
is a maximal chain of the same length.
Therefore $|\Phi(A)|=|A|$.

It follows that $\Phi$ induces, for each $r\ge 0$, a bijection between the subgroups of $G_{1}$
of order $p^{r}$ and the subgroups of $G_{2}$ of order $p^{r}$.
Hence $a_{p^{r}}(G_{1})=a_{p^{r}}(G_{2})$ for every $r$, and therefore
$\zeta_{G_{1}}(s)=\zeta_{G_{2}}(s)$.
\end{proof}

\end{document}
